\chapter{Implementación}
\label{cap:implementacion}

Este capítulo describe las decisiones de implementación más significativas del sistema SGROAS, con énfasis en la estrategia híbrida de acceso a datos y los patrones de código más relevantes.

\section{Configuración general}
\label{sec:configuracion}

El sistema se compila con Java 21 LTS y Spring Boot 3.5.x. La configuración principal se define en \texttt{application.properties} con valores por defecto para desarrollo y variables de entorno para producción (patrón 12-factor). Las dependencias principales se administran con Maven (\texttt{pom.xml}).

\section{Aislamiento de capas}
\label{sec:aislamiento}

El código fuente se organiza en paquetes que reflejan la arquitectura limpia \cite{martin2017}:

\begin{itemize}
  \item \texttt{controller}: recibe peticiones HTTP, valida entrada con DTOs, delega a servicios.
  \item \texttt{service}: lógica de negocio, orquestación de repositorios, transacciones.
  \item \texttt{repository}: acceso a datos mediante Spring Data JPA.
  \item \texttt{entity}: mapeo ORM con anotaciones JPA y \texttt{@NamedStoredProcedureQuery}.
  \item \texttt{dto}: contratos de entrada/salida como Java Records.
  \item \texttt{config}: configuración de seguridad, caché, CORS y manejo de excepciones.
\end{itemize}

\section{Manejo de errores globales}
\label{sec:errores}

El manejo de errores se centraliza en \texttt{GlobalExceptionHandler} (\texttt{@RestControllerAdvice}), que convierte excepciones en respuestas Problem Details (RFC 7807) \cite{rfc7807} con los campos \texttt{type}, \texttt{title}, \texttt{status}, \texttt{detail} e \texttt{instance}. Esto garantiza consistencia en todas las respuestas de error de la API.

\section{Esquema de caching}
\label{sec:caching}

La caché se implementa con Spring Cache + Redis \cite{redis}. Los listados paginados de conductores, vehículos, rutas y asignaciones se almacenan en Redis con TTL configurable. La invalidación se realiza automáticamente al crear, actualizar o eliminar registros mediante \texttt{@CacheEvict} en los métodos de servicio correspondientes.

\section{Mecanismo de seguridad}
\label{sec:seguridad-impl}

La pila de seguridad se compone de:

\begin{enumerate}
  \item \textbf{Filtro JWT} (\texttt{JwtAuthenticationFilter}): intercepta cada petición, extrae el token de la cookie HttpOnly, valida la firma y carga el usuario en el contexto de seguridad.
  \item \textbf{Autorización por rol}: \texttt{SecurityConfig} define reglas de acceso por endpoint según el rol (ADMIN, COORDINADOR, SEGURIDAD).
  \item \textbf{CORS}: configurado por origen específico en producción (\texttt{cors-spec} \cite{cors-spec}).
  \item \textbf{Rate limiting}: control de intentos de login mediante Redis (6 intentos, respuesta 429).
  \item \textbf{BCrypt}: contraseñas cifradas con factor de trabajo $\geq$ 10.
\end{enumerate}

\section{Estrategia híbrida de acceso a datos}
\label{sec:acceso-datos-impl}

La estrategia híbrida documentada en ADR-006 (Sección~\ref{sec:adr006}) se implementa con dos patrones claros:

\subsection{CRUD con Spring Data JPA}

Las operaciones CRUD elementales utilizan repositorios que extienden \texttt{JpaRepository}. El siguiente listado muestra el repositorio de conductores con consultas derivadas de nombres de métodos:

\begin{lstlisting}[language=Java, caption={Driver repository with derived queries.}, label={lst:conductor-repo}]
public interface DriverRepository extends JpaRepository<Driver, Long> {

    Page<Driver> findByActivoTrue(Pageable pageable);

    @Query("""
            SELECT c FROM Driver c
            WHERE c.activo = true
              AND (:search IS NULL
                   OR LOWER(c.nombres) LIKE LOWER(CONCAT('%', :search, '%'))
                   OR LOWER(c.apellidos) LIKE LOWER(CONCAT('%', :search, '%'))
                   OR c.cedula LIKE CONCAT('%', :search, '%')
                   OR LOWER(c.numeroLicencia) LIKE LOWER(CONCAT('%', :search, '%')))
            """)
    Page<Driver> searchActive(@Param("search") String search,
                              Pageable pageable);
}
\end{lstlisting}

Los métodos \texttt{findByActivoTrue} y \texttt{searchActive} utilizan consultas derivadas y JPQL estática respectivamente, sin concatenación de SQL dinámico.

\subsection{Stored procedures con @NamedStoredProcedureQuery}

Las operaciones de reportes y agregaciones se encapsulan en stored procedures de PostgreSQL. El siguiente listado muestra la declaración en la entidad \texttt{Incidente} y la invocación desde el repositorio:

\begin{lstlisting}[language=Java, caption={@NamedStoredProcedureQuery declaration in Incident.}, label={lst:incidente-sp}]
@NamedStoredProcedureQuery(
    name = "Incident.incidentesPorGravedad",
    procedureName = "sp_incidentes_por_gravedad",
    parameters = {
        @StoredProcedureParameter(
            mode = ParameterMode.IN,
            name = "p_tipo", type = String.class),
        @StoredProcedureParameter(
            mode = ParameterMode.REF_CURSOR,
            name = "cur", type = Class.class)
    })
public class Incident { ... }
\end{lstlisting}

\begin{lstlisting}[language=Java, caption={SP invocation from the repository.}, label={lst:incidente-repo}]
public interface IncidentRepository
        extends JpaRepository<Incident, Long> {

    @Procedure(name = "Incident.incidentesPorGravedad")
    List<Object[]> incidentsBySeverity(
            @Param("p_tipo") String tipo);
}
\end{lstlisting}

El cursor REFCURSOR se lee dentro de la misma transacción JDBC, por lo que \texttt{ReportService} está anotado con \texttt{@Transactional}.

\subsection{Funciones SQL}

Las funciones SQL se invocan con \texttt{SELECT}. El siguiente listado muestra la función \texttt{fn\_total\_programaciones}:

\begin{lstlisting}[language=Java, caption={Native queries in ScheduleRepository.}, label={lst:fn-programaciones}]
public interface ScheduleRepository
        extends JpaRepository<Schedule, Integer> {

    @Query(value = "SELECT estado AS clave, COUNT(*) AS total "
                 + "FROM programacion GROUP BY estado "
                 + "ORDER BY total DESC",
           nativeQuery = true)
    List<CountProjection> countByStatus();

    @Query(value = """
            SELECT TO_CHAR(fecha, 'YYYY-MM') AS clave, COUNT(*) AS total
            FROM programacion
            WHERE fecha >= CURRENT_DATE - INTERVAL '180 days'
            GROUP BY TO_CHAR(fecha, 'YYYY-MM')
            ORDER BY clave
            """, nativeQuery = true)
    List<CountProjection> countByMonth();
}
\end{lstlisting}

\section{Estrategia de seguridad estática}
\label{sec:seguridad-estatica}

El script \texttt{scripts/audit-sql-dynamic.sh} se ejecuta en el pipeline de CI y verifica:

\begin{itemize}
  \item Ausencia de \texttt{EXECUTE IMMEDIATE}, \texttt{sp\_executesql} o equivalentes en \texttt{db/procs/*.sql}.
  \item Ausencia de concatenación con \texttt{||} en consultas SQL dentro de \texttt{db/procs/}.
  \item Ausencia de \texttt{@Query} o \texttt{createNativeQuery} con concatenación en \texttt{src/main/}.
  \item Ausencia de \texttt{ddl-auto=update} en la configuración.
\end{itemize}

El análisis estático con SpotBugs y Find-Sec-Bugs complementa esta verificación, enfocándose en la regla \texttt{SQL\_NONCONSTANT\_STRING\_PASSED\_TO\_EXECUTE}.

\section{Pipeline de CI}
\label{sec:ci}

El pipeline de GitHub Actions (\texttt{.github/workflows/ci.yml}) ejecuta las siguientes etapas:

\begin{enumerate}
  \item Build y pruebas (\texttt{./mvnw verify}).
  \item Reporte JaCoCo con umbral mínimo de 70\,\%.
  \item Auditoría SQL dinámico (\texttt{audit-sql-dynamic.sh}).
  \item Validación de trazabilidad (\texttt{validate-traceability.sh}).
  \item Análisis estático (SpotBugs).
\end{enumerate}

El pipeline muestra al menos tres ejecuciones consecutivas exitosas con badge passing en el README.
